\section{Introduction}\label{sec:intro}

The gradient-flow formalism as introduced by
L\"uscher~\cite{Luscher:2010iy} and further formalized by L\"uscher and
Weisz~\cite{Luscher:2011bx} has proven useful in lattice \qcd\ in many
respects. One of its main virtues is that composite operators at finite
flow time $t$ do not require ultra-violet (\abbrev{UV}) renormalization
beyond the one of the involved parameters and fields. This means that
the operators do not mix under renormalization-group running, which
makes it particularly simple to combine results from different
regularization schemes. This feature opens promising prospects for a
cross-fertilization of lattice and perturbative calculations, such as a
possible lattice determination of $\alpha_s(M_Z)$, for
example\,\cite{Harlander:2016vzb}.

A particularly powerful way to exhibit this possible interplay is
obtained by considering the expansion of composite operators in the
limit of small flow time, which expresses flowed operators in terms of
\qcd\ operators at $t=0$, with $t$-dependent Wilson
coefficients~\cite{Luscher:2011bx}. This method has been used by Makino
and Suzuki~\cite{Suzuki:2013gza,Makino:2014taa} to derive a
regularization-independent formula for the energy-momentum tensor (\emt)
$T_{\mu\nu}$ which has already led to promising results (see, e.g.,
\citeres{Asakawa:2013laa,Kitazawa:2016dsl,Taniguchi:2016ofw,Kitazawa:2017qab,
 Yanagihara:2018qqg}).

The universal Wilson coefficients that occur in the formula of
\citere{Makino:2014taa} for the \abbrev{EMT} have been calculated
through next-to-leading order (\nlo) in perturbation
theory~\cite{Suzuki:2013gza,Makino:2014taa}. This corresponds to a
\one-loop calculation in the sense that it involves integrals over a
single $D$-dimensional momentum. In this paper, we will carry this
calculation to the next perturbative order.\footnote{Note that we work
  in the limit of infinite volume. The inclusion of finite-volume
  effects requires different techniques, such as Numerical Stochastic
  Perturbation Theory, see \citere{DallaBrida:2017tru}.}  It is
important to note at this point that the integrals which occur in the
gradient-flow formalism are of a more general type than in regular
\qcd. They involve additional exponential factors which depend on loop
and external momenta, as well as on flow-time variables, some of which
are also integrated over. Nevertheless, the first two-loop result was
already obtained in \citere{Luscher:2010iy}, even in analytic form.  The
extension to the three-loop level required significant aid from computer
algebra and numerical tools\,\cite{Harlander:2016vzb}.  From the
quantum-field theoretical point of view, it closely followed the steps
of \citere{Luscher:2010iy} by directly expressing the Green's functions
in terms of integrals with the help of Wick's theorem. The integrals
themselves were evaluated using
sector-decomposition\,\cite{Binoth:2003ak,Smirnov:2013eza} in order to
isolate the poles in $D-4$, whose coefficients were determined using
high-precision numerical methods~\cite{genz,holodborodko}.

In the current calculation, we apply a completely independent setup.  On
the one hand, it applies the gradient-flow formalism described in terms
of a five-dimensional quantum field theory\,\cite{Luscher:2011bx}, which
leads to well-defined, albeit non-standard Feynman rules. On the other
hand, rather than evaluating the resulting integrals numerically, we
express them in terms of master integrals using the integration-by-parts
method of Chetyrkin and Tkachov\,\cite{Chetyrkin:1981qh}. This reduces
the \nlo\ calculation of the Wilson coefficients of the \emt\ to a
single \one-loop integral without flow-time integration. The
next-to-next-to-leading order (\nnlo) calculation leads to four
\two-loop master integrals without flow-time integration, and two
\two-loop master integrals with a single flow-time integration. All
master integrals can be calculated analytically by standard means for
general values of $D$, the number of space-time dimensions.

By suitable renormalization, the Wilson coefficients of the \emt\ can be
defined in such a way that they are formally renormalization-scale
independent. For a fixed-order perturbative result, this means that the
renormalization-scale dependence is formally of higher order. This
allows one to estimate the perturbative uncertainty on the Wilson
coefficients through their residual dependence on the renormalization
scale $\mu$ around a particular ``central'' value. Based on the form of
the analytical result, we argue for a specific choice of this central
value. Our numerical study shows that the higher-order terms indeed lead
to an appreciable reduction of the $\mu$-variation. However, by
comparison of the successive higher-order terms, it appears that the
uncertainty estimate from a variation within $\mu\in[\mu_0/2,2\mu_0]$,
as it is common practice in regular perturbative \qcd\ calculation,
might be too optimistic.

While we consider the \nnlo\ expressions for the Wilson coefficients of
the \emt\ as our main result, our calculation allows us to obtain a
number of additional results that might be useful in a broader
context. Among these are the flowed quark-field renormalization constant
$Z_\chi$, and the matrix of anomalous dimensions for the set of
operators which form the energy-momentum tensor in regular (non-flowed)
\qcd\ through \nnlo.

The remainder of the paper is structured as follows. After briefly
introducing the perturbative gradient-flow formalism in order to define
our notation in \sct{sec:pert}, we outline the approach of
\citeres{Suzuki:2013gza,Makino:2014taa} for using this formalism to
define the \emt\ in \sct{sec:emt}. Technical details of our calculation
are described in \sct{sec:calc}. Section\,\ref{sec:coefs} contains our
main result, the Wilson coefficients through \nnlo\ \qcd\ in the
\msbar\ scheme. As pointed out in \citere{Makino:2014taa}, the trace
anomaly of the \emt\ allows for a welcome check of the calculation; we
briefly describe the derivation of the resulting relations among the
coefficient functions in \sct{sec:traceanom}. Finally, in \sct{sec:zij}, we
use the finiteness condition of the flowed operators in order to derive
the anomalous-dimension matrix for the set of operators which occur in
the \emt\ in regular \qcd. Section\,\ref{sec:conclusions} presents our
conclusions.
